\input{../../../stock/en-tete_v6.tex}

\newcommand{\titrechapitre}{titre}
\newcommand{\numerochapitre}{numéro}

\renewcommand{\headrulewidth}{0.4pt}
\renewcommand{\footrulewidth}{0.4pt}
\fancyfoot[L]{Notes de Raphaël JONTEF}
\fancyfoot[C]{\thepage}
\fancyfoot[R]{Enseirb-Matmeca}
\fancyhead[L]{Cours de Mathématiques}
\fancyhead[R]{\titrechapitre}

\begin{document}
\input{../../../stock/commands.tex}

% Page de titre custom
\setlength{\headheight}{15pt}
\begin{titlepage}
    \centering
    \vspace*{3cm}
    {\huge Chapitre \numerochapitre\par}
    {\Huge \bfseries\titrechapitre\par}
    \vspace{2cm}
    {\Large Notes rédigées par Raphaël JONTEF\par}
{\Large Avec l'aide de \par}
    \vspace{0.5cm}
    {\normalsize Cours enseigné par \par}
    \vspace{4cm}
    {\small Enseirb-Matmeca\par}
    {\small filière informatique\par}
    \vfill
    {\large \today\par}
\end{titlepage}

% plan du cours
\tableofcontents
\newpage


% -----------Corps du document--------------

Dans le second cas, si la complexité en temps de l'appel \code{opérateur(T , i , j)} est en $\Theta\big(\min(k-i,j-k)\big)$, On note $n = j - i + 1$ la longueur du tableau $T[i..j]$.

Si $k$ désigne la position relative de l’opérateur renvoyée par 
\code{indiceOpérateur} (c’est-à-dire $k = p-i$ si l’opérateur est en $p$), 
alors la complexité $\#(n)$ de \code{evalInfixe} vérifie~:

\[
\#(n)=
\begin{cases}
\Theta(1) & \text{si } n \leq 2\\
\underbrace{\#(k)}_{\text{appel réc. à gauche}} + \underbrace{\#(n-k-1)}_{\text{appel réc. à droite}} + \underbrace{\Theta\bigl(\min(k, n-k-1)\bigr)}_{\text{hypothèse}}
& \text{sinon}
\end{cases}
\]
où $k$ parcourt l’intervalle $1 \le k \le n-2$.

Cette borne vient du fait qu’un opérateur infixe doit laisser au moins un 
élément à gauche et un élément à droite : les sous-tableaux $T[i..p-1]$ 
et $T[p+1..j]$ ont donc des tailles
\[
k = p-i \quad\text{et}\quad n-k-1 = j-p
\]
qui doivent être toutes deux $\ge 1$, d’où la contrainte 
$1 \le k \le n-2$.

D'après la proposition supposée~    :

\[
\forall n\ge 1,\qquad
\#(n) \le 1 + n\log n
\]

On en déduit alors que la complexité temporelle de \code{evalInfixe} est en~:
\[
\#(n) = O(n \log n)
\]




\end{document}